
\textbf{Datasets.}

\begin{itemize}
\item \textit{Small CNN Zoo (CIFAR-10-GS)} \citep{Unterthineretal2020}: A collection of small CNNs with full training trajectories (epochs 0-50). The architecture is dynamically inferred from checkpoint state dict shapes (channel widths vary across models; the typical configuration observed is 3 convolutional layers followed by 2 fully-connected layers, with GELU activations and no padding). H-E1 used 200 CNN Zoo checkpoints (2,000 permutation runs); H-M1 used 100; H-M2 used 1,000 models x 50 epochs = 50,000 checkpoints.
\end{itemize}

\begin{itemize}
\item \textit{Small Transformer Zoo (MNIST)} [Tran-Viet et al., 2024]: A subset of the 125,000 checkpoint transformer zoo. The architectures used are 2-block ViT-style models with separate Q/K/V projections, no LayerNorm, and a 2-layer MLP classifier. H-E1 used 250 Transformer Zoo checkpoints (2,500 permutation runs); H-M1 used 100.
\end{itemize}

\textbf{Evaluation metrics.}

\begin{table}[htbp]
\centering
\begin{tabular}{lllll}
\toprule
Metric & Hypothesis & Threshold &  &  \\
\midrule
Mean & Delta acc &  & H-E1 & < 0.001 \\
Orbit-PE success rate & H-E1 & = 1.0 &  &  \\
overhead_ratio_mean & H-M1 & <= 1.20x &  &  \\
HAS_ARCH_BRANCHES & H-M1 & False &  &  \\
Var_perm/(Var_perm+Var_GL) & H-M2 & > 0.60 &  &  \\
\bottomrule
\end{tabular}
\end{table}

\textbf{Implementation.} All experiments ran on CPU for orbit-PE and variance computations. The symmetry validation experiment (H-E1) was executed on a single GPU (NVIDIA H100 NVL, CUDA_VISIBLE_DEVICES=2) in approximately 16.2 seconds. Orbit-PE uses d_pe = 64. Random seeds were fixed (seed = 42 for checkpoint sampling; seeds 0-9 for permutation sampling in H-E1). Code: \texttt{h-e1/code/orbit_pe.py} (dispatch table), \texttt{h-m1/code/orbit_pe_computer.py} (timing), \texttt{h-m2/code/orbit_projector.py} (variance decomposition), \texttt{h-m2/code/evaluate.py} (zoo-scale pipeline).

\textbf{Baseline comparisons not conducted.} Formal comparison of tau values against published NFN or SANE baselines was not conducted in this work. H-M3 (which would have trained SANE+orbit-PE and measured tau_retention) was blocked by the H-M2 gate failure. Published baseline values referenced in the text (NFN: tau = 0.934, 0.931; Transformer-NFN: tau = 0.905-0.910; SANE: linear probe 0.978-0.991) are taken directly from the respective papers and are not reproduced here.

\textbf{SVHN Zoo data unavailability.} The pre-specified cross-dataset stability check for H-M2 (|CIFAR ratio - SVHN ratio| < 0.10) was not performed because SVHN Zoo data was unavailable. The H-M2 conclusions are based solely on the CIFAR-10-GS subset.

